%%ha_lyman_td
Dans un premier temps, l'algorithme trie les valeurs du TD pour garder
uniquement les points o\`u le capteur est proche de l'�quilibre.
La s\'election se fait suivant deux crit\`eres~:\\
\indent - valeur de la tension de l'\'el\'ement Peltier proche de
la valeur \`a l'\'equilibre (compris entre $V_{pelI}$ et $V_{pelS}$)\\
\indent - $Td \leq Ts$ (suppression des d\'epassements)\\
\indent - $dp \geq 12.$ (vitesse air suffisante pour ventilater les capteurs)\\
\\
La s�rie d'humidit\'e absolue de r\'ef\'erence est calcul\'ee \`a partir
de la s\'erie des valeurs correctes du TD.\\
Une r\'egression lin\'eaire est r\'ealis\'ee entre, les
points d'humidit\'e absolue de r\'ef\'erence et, les points
correspondants du Lyman-$\alpha$.
Le Lyman-$\alpha$ est ensuite calibr\'e avec le r\'esultat de cette
r\'egression.

